% ============================================================
%  Поглавље 1: Увод
% ============================================================
\section{Увод}
\label{sec:uvod}

Рачунарство је суочено са класом проблема које класични рачунари
не могу да решавају ефикасно за велике улазе, такозваним НП-тешким
проблемима. Ови проблеми јављају се у логистици, криптографији,
машинском учењу и верификацији софтвера, па је њихово ефикасно решавање
од велике практичне важности.


Квантно рачунарство нуди нову рачунску парадигму засновану на принципима
квантне механике, суперпозицији и спрегнутости, која потенцијално
омогућује истовремено претраживање експоненцијално великог простора стања.
Квантна надмоћ над класичним рачунарима за НП-тешке проблеме
теоријски није доказана, али хибридни квантно-класични алгоритми показују
значајан практичан потенцијал.

У овом раду ћемо користи проблем \textit{MaxCut} као водећи пример кроз који
ћемо покушати да приближимо могуће предности једне гране квантног рачунарства.
Представићемо модел \textbf{квантног жарења (Quantum Annealing)} и одговарајући математички формализам 
 \textbf{квадратне неограничене бинарне оптимизације (QUBO quadratic unconstrained binary optimisation)}
и на њему моделовати проблем \textit{MaxCut}. Такође ћемо исти проблем моделовати
коришћењем \textit{Z3} решавача и указати на везу са \textit{QUBO}.